%!TEX program = xelatex
\documentclass{beamer}

\usepackage[UTF8]{ctex} 

\usetheme{metropolis}  
\title{BNU Beamer模板}
\subtitle{基于Metropolis设计的BNU Beamer模板}
\date{\today}
\author{Halve Luve}
\institute{北京师范大学人工智能学院}
\titlegraphic{
    \flushright\includegraphics[height=1cm]{BNU/BNU_name.pdf}
}
\addbibresource{ref.bib}

\begin{document}
    \maketitle
    \tableofcontents

    \section{部分一}
    \begin{frame}{无序列表}
        以下是无序列表的例子:
        \begin{itemize}
            \item 第一点
            \item 第二点
        \end{itemize}
    \end{frame}

    \begin{frame}{分列展示}
        \begin{columns}
            \column{.5\textwidth}
            这是第一列
            
            \column{.5\textwidth}
            这是第二列
        \end{columns}
    \end{frame}

    \section{部分二}

    \begin{frame}{表格}
        \begin{table}
            \centering
            \begin{tabular}{ccc}
                \hline
                \textbf{标签1} & \textbf{标签2} & \textbf{标签3} \\
                \hline 
                值1 & 值2 & 值3 \\ 
                值4 & 值5 & 值6 \\ 
                值7 & 值8 & 值9 \\ 
                \hline
            \end{tabular}
            \caption{表1的题注}
            \label{tab:table1}
        \end{table}
    \end{frame}

    \begin{frame}{图像}
        \begin{figure}[H]
            \centering
            \includegraphics[width=.5\textwidth]{example-image-a}
            \caption{图1的题注}
            \label{fig:figure1}
        \end{figure}
    \end{frame}

    \begin{frame}{引用}
        这里引用了一张图 (图\ref{fig:figure1})\\
        这里引用了一个表 (表\ref{tab:table1})\\
        这里引用了一篇文章 \cite{vaswani2017attention}\\
        % \footnote{不推荐这种引用方式，因为\textbf{printbibliography}在这个主题下看起来不太正常。未来可能会修复。也可通过PR或者Issues提供建议 }
        这里以脚注形式引用了一篇文章 \footfullcite{vaswani2017attention}
    \end{frame}
    
    \lecture{Test lecture one}{lone}
    \begin{frame}{定理}
        \begin{theorem}
            这是定理1。
        \end{theorem}
        \begin{definition}
            这是一个定义。
        \end{definition}
        \begin{corollary}
            这是一个推论。
        \end{corollary}
    \end{frame}

    \lecture{Test lecture two}{ltwo}
    \begin{frame}{公式}
        \begin{theorem}
            这是定理2。
        \end{theorem}
        这是公式块：
        \begin{equation}
            H(x)=-\sum\limits_i^Np(x_i)\log p(x_i).
        \end{equation}
        这是内联公式：$H(x)=-\sum\limits_i^Np(x_i)\log p(x_i)$。
    \end{frame}

    \section{参考文献}

    \begin{frame}[allowframebreaks]{参考文献}
        % \tiny\bibliography{references}
        % \bibliographystyle{apalike}
        \printbibliography[heading=none]
    \end{frame}
\end{document}